\chapter{Teoría espectral}


\noindent 
Sea $F: X \to X$ una aplicación lineal y continua en un espacio de Banach $X$, entonces hemos visto que si $F$ es biyectiva, entonces $F^{-1}$ es continua. Sea 
$$F := T - \lambda Id$$
con $T \in BL(X)$, preguntémonos cuándo $F$ es biyectiva. \\

\begin{definicion}[Resolvente]
    Sea $T \in BL(X)$, definimos 
    $$\rho(T) := \{ \lambda \in \mathbb{R} \mid T - \lambda Id \text{ es biyectiva} \}$$
    como el \tbf{resolvente}.
\end{definicion}
\begin{definicion}[Espectro]
    Sea $T \in BL(X)$, definimos 
    $$\sigma(T) := \{ \lambda \in \mathbb{R} \mid T - \lambda Id \text{ no es biyectiva} \}=\R\setminus \rho(T)$$ 
    como el \tbf{espectro}. \\
\end{definicion}


\observacion{
    Si $\lambda \in \sigma(T)$, entonces $T - \lambda Id$ o no es inyectiva o no es sobreyectiva.
}

\begin{definicion}
    Sea $T\in BL(X)$, entonces $\lambda \in \sigma(T)$ se dice \textbf{valor propio} de $T$ si $T - \lambda Id$ no es inyectiva. Definimos 
    $$e(T) := \{ \lambda \in \sigma(T) \mid \lambda \text{ es valor propio} \} \subseteq \sigma(T)$$
    como el \tbf{conjunto de los valores propios}.
\end{definicion}

\observacion{
    Sea $\lambda \in e(T)$, entonces 
    $$\exists x_1 \neq x_2 \in X : (T - \lambda Id)(x_1) = (T - \lambda Id)(x_2)$$
    Sea $y:=x_1-x_2\neq 0$, entonces
    \begin{equation*}
        (T - \lambda Id)(y) = 0 \iff T(y) = \lambda y
    \end{equation*}
    De lo cual se deduce que 
    \begin{itemize}
        \item $\lambda \in e(T) \iff \exists y \neq 0 : T(y) = \lambda y$ (ecuación de valores propios).
        \item $\lambda \in e(T) \iff \text{Ker}(T - \lambda Id) \neq \{0\}$.
    \end{itemize}
    En este caso, $y$ se llama \textbf{vector propio} de $T$ respecto a $\lambda$ y $\text{Ker}(T - \lambda Id)$ se llama \textbf{espacio propio} de $\lambda$, el cual es un subespacio cerrado.
}


\begin{teo}
    Sea $X$ un espacio de Banach y sea $T \in BL(X, Y)$, entonces
    \begin{enumerate}
        \item Si $\lambda_1, \lambda_2 \in e(T)$ con $\lambda_1 \neq \lambda_2\implies \text{Ker}(T - \lambda_1 I) \cap \text{Ker}(T - \lambda_2 I) = \{0\}$.
        \item Si $\{v_1, \dots, v_m\}$ son vectores propios correspondientes a valores propios distintos \mbox{$\lambda_1, \dots, \lambda_m \in e(T)$}, entonces $\{v_1, \dots, v_m\}$ es un conjunto de vectores linealmente independientes.
    \end{enumerate}
    %//TODO: leer demostración y ver si se copio bien
\begin{proof}
    \begin{description}
        \item[$1)$] Sea $\bar{x} \in \text{Ker}(T - \lambda_1 I) \cap \text{Ker}(T - \lambda_2 I)$, entonces $T(\bar{x}) = \lambda_1 \bar{x} = \lambda_2 \bar{x} \iff \underbrace{(\lambda_1 - \lambda_2)}_{\neq 0} \bar{x} = 0 \implies \bar{x} = 0$.
        \item[$2)$] Sean $E_1, \dots, E_m$ los espacios propios de $\lambda_1, \dots, \lambda_m$ respectivamente. Sean $\{x_1, \dots, x_m\}$ tales que $x_i \in E_i \quad \forall i = 1, \dots, m$. Entonces $\{x_i\}$ es linealmente independiente.

        Supongamos $k > 1$ y demostremos que $\{x_1, \dots, x_k\}$ es linealmente independiente sabiendo que $\{x_1, \dots, x_{k-1}\}$ lo es.

        Sea $\alpha_1 x_1 + \dots + \alpha_k x_k = 0 \implies 0 = \alpha_1 T(x_1) + \dots + \alpha_k T(x_k) \iff 0 = \alpha_1 \lambda_1 x_1 + \dots + \alpha_k \lambda_k x_k$.

        Supongamos por el contrario (P.A.) que $\alpha_k = 1 \implies \lambda_k x_k = -\alpha_1 \lambda_1 x_1 - \dots - \alpha_{k-1} \lambda_{k-1} x_{k-1}$.

        Pero $x_k = -\alpha_1 x_1 - \dots - \alpha_{k-1} x_{k-1}$, y por lo tanto:
        $$ -\lambda_k (\alpha_1 x_1 + \dots + \alpha_{k-1} x_{k-1}) = -\alpha_1 \lambda_1 x_1 - \dots - \alpha_{k-1} \lambda_{k-1} x_{k-1}, $$ de lo cual:
        $$\alpha_1 (\lambda_1 - \lambda_k) x_1 + \dots + \alpha_{k-1} (\lambda_1 - \lambda_{k-1}) x_{k-1} = 0$$
        Como $\{x_1, \dots, x_{k-1}\}$ son linealmente independientes $\implies$
        $$\alpha_i \underbrace{(\lambda_i - \lambda_k)}_{\neq 0} = 0 \quad \forall i \implies \alpha_1 = \dots = \alpha_{k-1} = 0 \implies x_k = 0 \text{ vector propio: ¡ABSURDO!}$$
    \end{description}
\end{proof}
\end{teo}

\begin{teo}
    Sea $X$ un espacio de Banach y sea $T \in BL(X, Y)$, entonces
    \begin{enumerate}
        \item $\rho(T)$ es abierto.
        \item $\sigma(T) \subseteq [-\|T\|, \|T\|]$.
    \end{enumerate}
    De lo cual $\sigma(T)$ es compacto.
\begin{proof}
    \begin{description}
        \item[$1)$] Sea $\alpha_0 \in \rho(T)$, es decir, $T - \alpha_0 Id$ es biyectiva. Demostremos que si $\alpha$ está cerca de $\alpha_0$, entonces $T - \alpha Id$ es biyectiva.

        \noindent Sea $f \in X$, consideremos $T(u) - \alpha u = f$. Si $\alpha \in \rho(T) \implies \exists! u : T(u) - \alpha u = f \iff$
        \[ \iff T(u) - \alpha u - \alpha_0 u = f - \alpha_0 u \iff T(u) - \alpha_0 u = f - \alpha_0 u + \alpha u. \quad \text{Aplicamos } (T - \alpha_0 Id)^{-1} \implies \]
        \[ u = (T - \alpha_0 Id)^{-1}(f - \alpha_0 u + \alpha u) =: G(u). \quad \text{Por lo tanto:} \]
        $\alpha \in \rho(T) \iff G(u) = u$. Si $G$ es una contracción, hemos terminado:
        \[ \|G(u) - G(v)\| = \|(T - \alpha_0 Id)^{-1}(f - \alpha)(u - v)\| \leq \|(T - \alpha_0 Id)^{-1}\| \cdot |f - \alpha| \cdot \|u - v\| \]
        Sea $\delta > 0 : \delta < \frac{1}{\|(T - \alpha_0 Id)^{-1}\|}$ para $|\alpha - \alpha_0| < \delta \implies G$ es una contracción.
        \item[$2)$]Si $\alpha \in \mathbb{R} : |\alpha| > \|T\|$, mostremos que $\alpha \in \rho(T)$:

        \noindent $T(u) - \alpha(u) = f$, mostremos que $\forall f \in X \ \exists!$ solución de $T(u) - \alpha u = f$.
        \[ u := \frac{1}{\alpha}(T(u) - f) =: G(u) \implies \|G(u) - G(v)\| = \frac{1}{|\alpha|} \|T(u) - T(v)\| \leq \frac{\|T\|}{|\alpha|} \|u - v\| < 1 \text{ si } |\alpha| > \|T\|\]
    \end{description}
\end{proof}
\end{teo}

\begin{ejemplo}
    \begin{enumerate}
        \item $T = Id$. Resolvamos $(T - \lambda Id)(x) = 0$ con $x\neq 0$ para buscar los valores propios 
        \begin{equation*}
                Id(x) - \lambda Id(x) = 0\iff x - \lambda x = 0 \iff (1 - \lambda)\underbrace{x}_{\neq 0} = 0 \iff \lambda = 1
        \end{equation*}
        por lo que $\lambda = 1$ es un valor propio con espacio propio $X$. Si $\lambda \neq 1$, entonces 
        $$Id(x) - \lambda Id(x) = f \implies x(1 - \lambda) = f \implies x = \frac{1}{1 - \lambda} f$$ 
        donde entonces hemos demostrado que $T$ es biyectiva para $\lambda \neq 1$, pues justo ahora se demostró la sobreyectividad y $T$ es inyectiva para $\lambda \neq 1$, por lo que $\rho(T) = \mathbb{R} \setminus \{1\}$.
        \item En este ejemplo trabajaremos con $X = \ell_p$ y el operador
        \Func{T}{\ell_p}{\ell_p}{x}{\left(x_1, \frac{x_2}{2}, \dots, \frac{x_n}{n}, \dots\right)}

        Si planteamos la ecuación de valores propios $T(x) = \lambda x$ término a término para cada posición de la sucesión, obtenemos un sistema infinito de ecuaciones independientes. \\
        \noindent
        Vamos a resolver la ecuación para una posición genérica $n$. Si pasamos todo a un solo lado restando, podemos sacar factor común $x_n$
        $$\frac{x_n}{n} - \lambda x_n = 0 \iff \left(\frac{1}{n} - \lambda\right) x_n = 0$$
        donde tenemos que o bien $x_n = 0$ o bien $\lambda = \dfrac{1}{n}$, donde por elección de la componente supondremos que $x_n=\alpha\neq 0$, puesto que sabemos que debe haber al menos una componente no nula para que $x$ sea un vector propio. Tenemos entonces que el vector propio es 
        \begin{equation*}
            \lambda=\dfrac{1}{n}
        \end{equation*}
        Para el resto de componentes $m\in \N$ con $m\neq n$, se tiene que
        \begin{equation*}
            \frac{x_m}{m} - \lambda x_m = 0 \iff \left(\frac{1}{m} - \frac{1}{n}\right) x_m = 0 \iff x_m = 0
        \end{equation*}
        donde comprobamos efectivamente que $\lambda$ es valor propio de la siguiente manera 
        $$T(0, \dots, \alpha, \dots) = \left(0, \dots, \frac{1}{n}, \dots\right) = \frac{1}{n}(0, \dots, \alpha, \dots)$$
        Por tanto, aplicando esto a cada componente $n\in \N$, obtenemos que 
        $$e(T) = \left\{ 1, \frac{1}{2}, \frac{1}{3}, \dots, \frac{1}{n}, \dots \right\} = \left\{ \frac{1}{n} \;\middle|\; n \in \mathbb{N} \right\}$$
    \end{enumerate}
\end{ejemplo}

\section{Operadores sobre los complejos}

\noindent 
Sea $T: H \to H$ un operador sobre $H$, un $\mathbb{C}$-espacio de Hilbert, entonces:
\begin{enumerate}
    \item[i)] Si $T$ es autoadjunto: $(T(x), x) = (x, T(x)) = \overline{(T(x), x)} \quad \forall x \in H \implies (T(x), x) \in \mathbb{R}$. \\
\end{enumerate}

\noindent 
Sea $X$ un espacio de Banach sobre $\mathbb{C}$ y $T \in BL(X)$. Definimos:
\[ \rho(T) := \{ \lambda \in \mathbb{C} \mid T - \lambda Id \text{ es biyectiva} \}, \quad \sigma(T) := \mathbb{C} \setminus \rho(T) \]

\begin{prop}
    Sea $T\in BL(\mathcal{X})$ con $X$ un espacio de Banach sobre $\mathbb{C}$, entonces \begin{enumerate}
        \item $\rho(T)$ es abierto en $\mathbb{C}$.
        \item $\sigma(T) \subseteq \{ \lambda \in \mathbb{C} \mid |\lambda| \leq \|T\| \}$, de lo cual $\sigma(T)$ es compacto.
    \end{enumerate}
\end{prop}

\begin{definicion}
    Definimos el \textbf{radio espectral} como  
    $$r(T) := \inf \{ r \geq 0 \mid \sigma(T) \subseteq B_r(0) \}$$
    y se llama \textbf{conjunto de valores propios} de $T$ a lo siguiente
    $$e(T) = \{ \lambda \in \sigma(T) \mid T - \lambda Id \text{ no es inyectiva} \}$$    
    Si $\lambda \in e(T)$, entonces $\dim(\text{Ker}(T - \lambda Id))$ se llama \textbf{multiplicidad geométrica} de $\lambda$.
\end{definicion}

\begin{teo}
    Sea $T \in BL(X)$, con $X$ espacio de Banach sobre $\mathbb{C}$, entonces 
    $$\sigma(T) \neq \emptyset$$
\end{teo}

\begin{ejemplo}
    Sea $X = \ell_2$, y sea el siguiente operador lineal
    \Func{T}{X}{X}{(x_1, \dots, x_m, \dots)}{(x_2, \dots, x_{m+1}, \dots)}
    con $\|T\| = 1$, entonces sabiendo que $\sigma(T) \subseteq \{ \lambda \in \mathbb{C} \mid |\lambda| \leq 1 \}$, calculemos $\sigma(T)$. Para ello resolvemos lo siguiente
    \begin{align*}
        T(x) = \lambda x &\iff
        \begin{cases} 
            x_2 = \lambda x_1 \\ 
            x_3 = \lambda x_2 \\ 
            \vdots \\ 
            x_{m+1} = \lambda x_m 
        \end{cases}
        \overset{x \neq 0}{\Longrightarrow}
        \begin{cases} 
            x_2 = \lambda x_1 \\ 
            x_3 = \lambda x_2 = \lambda^2 x_1 \\ 
            \vdots \\ 
            x_m = \lambda^{m-1} x_1 \\ 
            x_{m+1} = \lambda^m x_1 
        \end{cases}
    \end{align*}
    La sucesión $(x_1, \lambda x_1, \dots, \lambda^m x_1, \dots)$ satisface $T(x) = \lambda x$, pero:
    \begin{align*}
        x \in \ell_2 &\implies x_m \to 0 \implies \lambda^m x_1 \xrightarrow{m \to +\infty} 0 \iff |\lambda| < 1
    \end{align*}
    de lo cual definiendo este nuevo vector
    $$\tilde{x} = (1, \lambda, \dots, \lambda^m, \dots)$$
    sabemos que cumple $T(\tilde{x}) = \lambda \tilde{x}$, pues solo se ha tomado $x_1 = 1$ para que el vector propio no sea el vector nulo. Ahora para ver que efectivamente es vector propio, basta comprobar que $\tilde{x} \in \ell_2$ si $|\lambda| < 1$, lo cual es cierto, pues 
    $$\|\tilde{x}\|_{\ell_2}^2 = \sum_{m=0}^{\infty} |\lambda^{m}|^2 = \sum_{m=0}^{\infty} (|\lambda|^2)^m<\infty \iff |\lambda| < 1$$
    Por lo tanto, $\tilde{x}\in \ell_2$ si $|\lambda| < 1$ y como $T(\tilde{x}) = \lambda \tilde{x}$, entonces $\lambda$ es un valor propio si $|\lambda| < 1$. \\ \\
    Para $|\lambda| = 0$, tenemos que $\tilde{x} = (1, 0, \dots, 0)$ es un vector propio, entonces
    \begin{equation*}
        \{ |\lambda| < 1 \} \subseteq e(T) \subseteq \sigma(T) \subseteq \{ |\lambda| \leq 1 \}
    \end{equation*}
    Como $\sigma(T)$ es cerrado, se sigue que
    $$\overline{\{ \lambda \in \mathbb{C} \mid |\lambda| < 1 \}}=\{ \lambda \in \mathbb{C} \mid |\lambda| \leq 1 \} \subseteq \sigma(T)\implies\sigma(T) = \{ \lambda \in \mathbb{C} \mid |\lambda| \leq 1 \} $$
    donde hemos usado que un conjunto cerrado que contiene a otro conjunto $A$, debe contener a la frontera de $A$. \\ \\
    \noindent
    Calculemos ahora $T^*$, es decir, dicho operador que cumple
    \begin{equation*}
        ( x, T(y) ) = ( T^*(x), y ) \quad \forall x, y \in X
    \end{equation*}
    entonces
    \begin{align*}
        ( x, T(y) )_{\ell_2} &= \sum_{m=1}^{\infty} x_m \cdot y_{m+1} = \sum_{m=2}^{\infty} x_{m-1} y_m = \sum_{m=1}^{\infty} (T^*(x))_m y_m = ( T^*(x), y )_{\ell_2}
    \end{align*}
    de lo cual se deduce que
    \begin{equation*}        
        T^*(x_1, x_2, \dots, x_m, \dots) = (0, x_1, x_2, \dots, x_{m-1}, \dots)
    \end{equation*}
\end{ejemplo}

\begin{ejercicio}
    Demuestra lo siguiente
    \begin{equation*}
        \lambda \in \sigma(T) \iff \bar{\lambda} \in \sigma(T^*)
    \end{equation*}
    Si $\lambda \in \sigma(T)$, entonces $T - \lambda Id$ no es biyectiva, lo cual implica que $T - \lambda Id$ o no es inyectiva o no es sobreyectiva. En ambos casos, se deduce que 
    $$(T - \lambda Id)^*=T^*-\bar{\lambda} Id$$ 
    no es biyectiva, lo cual implica que $\bar{\lambda} \in \sigma(T^*)$. El recíproco se demuestra de la misma manera. \\
\end{ejercicio}

\begin{teo}
    Sea $T$ un operador compacto, entonces 
    $$\text{Ker}(T - \lambda Id)$$
    es de dimensión finita $\forall \lambda \neq 0$.
\begin{proof} % //TODO: leer demostración y ver si se copio bien
    Basta mostrar que la bola unitaria es compacta $B_1 \subseteq \text{Ker}(T - \lambda Id)$ es compacta
    \begin{align*}
        (x_n) \subseteq B_1 : \|x_n\| \leq 1 \quad : \quad (T - \lambda Id)(x_n) = 0 \quad \forall n \implies T(x_n) = \lambda x_n \quad \forall n
    \end{align*}


    \noindent $(x_n)$ es acotada y $T$ es compacto $\implies \exists$ subsucesión convergente $T(x_{n_k}) \to z \implies \lambda x_{n_k} \to z$

    \noindent $x_{n_k} \to z/\lambda =: x_0 \quad \text{con} \quad \|x_0\| = \lim_{k} \|x_{n_k}\| \leq 1 \implies x_0 \in B_1 \implies B_1 \text{ es compacta}$. \\
\end{proof}
\end{teo}


\begin{ejercicio}
    Si $X$ es de dimensión infinita y $T \in K(X)$, entonces 
    $$0 \in \sigma(T)$$
    Si por el contrario (P.A.) $T - 0 \cdot Id = T$ fuera biyectiva, entonces $T \circ T^{-1} = Id \in K(X)$ (ya que la composición de operadores compactos es compacta), ¡ABSURDO!.
\end{ejercicio}

\begin{teo}
    Sea $T \in K(X)$, entonces
    \begin{enumerate}
        \item[i)] $e(T)$ es a lo sumo numerable.
        \item[ii)] Si $e(T)$ no es finito, entonces $0$ es el único punto de acumulación.
    \end{enumerate}
\begin{proof}   % //TODO: leer demostración y ver si se copio bien
    \begin{description}
        \item[$i)$] Sea $E_m = \left\{ \lambda \in e(T) \mid |\lambda| \geq \frac{1}{m} \right\} \quad \forall m$ y supongamos por $(*)$ que $E_m$ es finito o vacío. Entonces:
        \[ e(T) = \bigcup_{m \geq 1} E_m \implies \text{unión numerable de conjuntos finitos} \implies e(T) \text{ es numerable}. \]

        \item[$ii)$] $e(T)$ es numerable $\implies$ posee al menos un punto de acumulación (Bolzano-Weierstrass). Mostremos que es $\tilde{\lambda} = 0$, de hecho:

        \noindent si $\tilde{\lambda} \neq 0$ fuera de acumulación, $\exists k > 0 : |\tilde{\lambda}| > \frac{1}{k}$, entonces $\exists I_\delta(\tilde{\lambda})$ de radio $\delta > 0 :$
        \[ |\lambda| > \frac{1}{k} \quad \forall \lambda \in I_\delta \quad \text{y} \quad e(T) \cap I_\delta \neq \emptyset \implies e(T) \cap I_\delta \text{ es a lo sumo finito por } (*) \]
        lo cual es un ¡ABSURDO! porque $\tilde{\lambda}$ es de acumulación.

        \noindent \textbf{AFIRMACIÓN (CLAIM).} $E_m = \left\{ \lambda \in e(T) \mid |\lambda| \geq \frac{1}{m} \right\}$ es finito $\forall m$.

        \noindent Supongamos por el contrario (P.A.) que existe un $\bar{m} \geq 1$ tal que existe una sucesión $\{\lambda_m\} \subseteq e(T)$ con $|\lambda_m| \geq \frac{1}{\bar{m}} \quad \forall m$, siendo los $\{\lambda_m\}$ distintos.

        \noindent Sea $(x_m)_m$ la sucesión de vectores propios: $T(x_m) = \lambda_m x_m \implies \{x_m\}$ es un sistema de vectores linealmente independientes.

        \noindent Consideremos $Y_m := \text{span}(x_1, \dots, x_m) \supsetneq Y_{m-1} \quad \forall m$, subespacios vectoriales cerrados y propios.

        \noindent Por el \textbf{lema de Riesz}, $\exists y_m \in Y_m : \|y_m\| = 1 \quad \text{y} \quad \|y_m - x\| \geq \frac{1}{2} \quad \forall x \in Y_{m-1}$. Se cumple:

        \noindent $(T - \lambda_m Id)(Y_m) \subseteq Y_{m-1} \quad \forall m$. De hecho, si $v \in Y_m : v = \alpha_1 x_1 + \dots + \alpha_m x_m$, entonces:
        \[ (T - \lambda_m Id)(\alpha_1 x_1 + \dots + \alpha_m x_m) = \alpha_1 (T - \lambda_m Id)(x_1) + \dots + \alpha_m \underbrace{(T - \lambda_m Id)(x_m)}_{0} = \]
        \[ = \alpha_1 (\lambda_1 - \lambda_m) x_1 + \dots + \alpha_{m-1} (\lambda_{m-1} - \lambda_m) x_{m-1} \in Y_{m-1}. \]

        \noindent Sea $v = y_m \in Y_m$, entonces: $(T - \lambda_m Id)(y_m) \in Y_{m-1}$.

        \noindent Además, $T(Y_k) \subseteq Y_{m-1} \quad \forall k \in \{1, \dots, m-1\} \implies y_k = \sum_{i=1}^{k} \beta_i x_i \in Y_k \implies T(y_k) = \sum_{i=1}^{k} \beta_i \lambda_i x_i \in Y_k$.

        \noindent Sean $z := T(y_k) - (T - \lambda_m Id)(y_m) \in Y_{m-1} \quad \forall k \in \{1, \dots, m-1\} \implies \frac{z}{\lambda_m} \in Y_{m-1}$.

        \noindent $T(y_k) - T(y_m) = z - \lambda_m y_m$
        \[ \|T(y_k) - T(y_m)\| = \|\lambda_m y_m - z\| = |\lambda_m| \cdot \left\| y_m - \frac{z}{\lambda_m} \right\| \xrightarrow{\text{Riesz}} \geq |\lambda_m| \cdot \frac{1}{2} \geq \frac{1}{\bar{m}} \cdot \frac{1}{2}. \quad \square \]
    \end{description}
\end{proof}
\end{teo}


\observacion{
    Si $T$ es compacto, $X$ es de Banach y  $e(T) = \{\lambda_m\}_m$ es numerable entonces 
    \begin{equation*}
            e(T) \subseteq [-\|T\|, \|T\|] \implies (\lambda_m) \text{ es acotada con } 0 \text{ como punto de acumulación} \implies \lambda_m \to 0 \\
    \end{equation*}
}

\begin{ejercicio}
    Sea el operador 
    \Func{T}{\ell_2}{\ell_2}{(x_1, x_2, \dots, x_m, \dots)}{(x_1 \frac{x_2}{2}, \dots, \frac{x_{m}}{m}, \dots)}
    Calcular $\sigma(T)$ y $e(T)$, y comprobar que $0 \in \sigma(T)$ pero $0 \notin e(T)$. \\ \\
    \noindent
    En un ejercicio anterior se calculó el conjunto de valores propios que es
    \begin{equation*}
        e(T) = \left\{ \frac{1}{n} \;\middle|\; n \in \mathbb{N} \right\}
    \end{equation*}
    Este operador $T$ es un operador compacto (podemos visualizarlo porque sus coeficientes diagonales $\frac{1}{n}$ tienden a cero, lo que significa que el operador ``aplasta'' las dimensiones infinitas de forma continua). Como es un operador compacto en un espacio de dimensión infinita como $\ell_2$, entonces $0 \in \sigma(T)$, y tenemos que el espectro viene dado como 
    \begin{equation*}
        \sigma(T) =\left\{0\right\} \cup \left\{ \frac{1}{n} \;\middle|\; n \in \mathbb{N} \right\}
    \end{equation*}

\end{ejercicio}

\begin{teo}[Teorema de la alternativa de Fredholm] \label{teo:fredholm}
    Sea $H$ un espacio de Hilbert y $T \in BL(H)$ compacto. Entonces:
    \begin{enumerate}
        \item[a)] si $\lambda \neq 0$, entonces $R(T - \lambda I)$ es cerrado y se cumple $R(T - \lambda I) = \text{Ker}(T^* - \bar{\lambda} I)^\perp$
        \item[b)] si $\lambda \neq 0$, entonces $\text{Ker}(T - \lambda I) = \{0\} \iff R(T - \lambda I) = H$, es decir:
        \[ T - \lambda I \text{ es inyectivo} \iff T - \lambda I \text{ es sobreyectivo} \]
        \item[c)] $\dim \text{Ker}(T - \lambda I) = \dim \text{Ker}(T^* - \bar{\lambda} I)$
    \end{enumerate}
\end{teo}

\noindent 
Veamos una aplicación. Supongamos que queremos resolver 

$$T(u) - \lambda u = f, \ \text{con } \lambda \neq 0, \ f \in H \text{ y } T \text{ compacto}$$
lo que buscaremos es las condiciones de solución para esta ecuación, por lo que primero agruparemos la incognita $u$ a un lado de la ecuación, y el término independiente $f$ al otro lado, quedando la ecuación de la siguiente forma
\begin{equation*}
    (T - \lambda Id)(u) = f
\end{equation*}
Preguntarse si esta ecuación es resoluble (es decir, si existe algún $u \in H$ que al aplicarle el operador nos dé $f$) es exactamente lo mismo que preguntarse si el vector $f$ pertenece a la imagen o rango del operador. En lenguaje matemático:
$$\text{La ecuación es resoluble} \iff f \in R(T - \lambda Id)=\{ f \in H \mid \exists u \in H \text{ tal que } (T - \lambda Id)(u) = f \}$$
Como bien indica el apartado $b)$ del Teorema \ref{teo:fredholm}, tenemos que $T-\lambda Id$ es inyectivo, o no, véamos los dos casos:
\begin{itemize}
    \item \tbf{Caso inyectivo:} si el operador es inyectivo, significa que su núcleo es trivial, es decir, $\text{Ker}(T - \lambda Id) = \{0\}$. Al ser inyectivo, el apartado $b)$ nos obliga a concluir que también es sobreyectivo ($R(T - \lambda Id) = H$), y por tanto significa que la imagen es todo el espacio $H$, por lo que para cualquier $f \in H$ existe solución.  \\ 
    \noindent
    Además, ser inyectivo significa que no hay dos caminos que lleven al mismo sitio, por lo que esa solución es única ($\exists!$).  Esto demuestra matemáticamente \underline{$\forall f \in H \ \exists!$ solución}.
    \item \tbf{Caso no inyectivo:} si el operador no es inyectivo, significa que la ecuación homogénea (cuando pones $f = 0$) tiene soluciones distintas de cero
    $$(T - \lambda Id)(u) = 0 \implies T(u) = \lambda u$$
    Es decir, existen vectores propios asociados al autovalor $\lambda$. Además, por la teoría general de operadores compactos (que vimos en el teorema anterior del claim), el espacio de soluciones de un autovalor distinto de cero tiene que ser de dimensión finita. \\ \\
    \noindent
    Esto demuestra matemáticamente que \underline{la ecuación con $f = 0$ tiene un número finito de soluciones} \underline{linealmente independientes}. \\
\end{itemize}
\noindent
Entonces $\forall f \in H \ \exists!$ solución, o bien la ecuación con $f = 0$ tiene un número finito de soluciones linealmente independientes. Por lo tanto:

\begin{equation*}
    T(u) - \lambda u = f \text{ es resoluble } (f \in R(T - \lambda Id)) \iff f \in \text{Ker}(T^* - \bar{\lambda} Id)^\perp \text{ y por (c) tiene dimensión } n \\
\end{equation*}

\noindent
Vamos a resumir en un resultado claro lo que tenemos hasta ahora.
\begin{coro}
    Dado un operador
    \Func{T}{H}{H}{u}{T(u)}
    compacto, lineal y continuo, con $\mathcal{H}$ un espacio de hilbert de dimensión infinita tenemos que 
    \begin{itemize}
        \item $0 \in \sigma(T)$ (sea valor propio o no).
        \item $\sigma(T) \setminus \{0\} = e(T) \setminus \{0\}$ (apartado $b)$ del Teorema \ref{teo:fredholm}).
        \item Se tiene una de las siguientes alternativas:
        \begin{itemize}
            \item $\sigma(T) = \{0\}$
            \item $\sigma(T) \setminus \{0\} = e(T) \setminus \{0\}$ es un conjunto finito.
            \item $\sigma(T) \setminus \{0\} = e(T) \setminus \{0\}$ es una sucesión $\lambda_n$ que tiene a $\{0\}$ como único punto de acumulación y $\lambda_n \to 0$.
        \end{itemize}
    \end{itemize}
\end{coro}

\observacion{
    Para el caso $X$ un espacio de Banach y no de Hilbert, tenemos que también se cumple para $T$ compacto, lineal y continuo, que
    \begin{equation*}
        \sigma(T) \setminus \{0\} = e(T) \setminus \{0\}
    \end{equation*}
}

\begin{ejercicio}
    Sea el operador
    \Func{T}{C^0([0,1])}{C^0([0,1])}{x}{T(x)}
    donde $x: [0,1] \to \mathbb{R}\in C^0([0,1])$ (función continua), y $T(x)(t) := t\cdot x(t)$. Se pide:
    \begin{itemize}
        \item[i)] Calcular $\sigma(T)$.
        \item[ii)] Demuestra si $T$ es compacto o no.
    \end{itemize}
    Veámos la resolución de cada apartado por separado.
    \begin{description}
        \item[$i)$] Tomemos una función arbitraria $f \in C^0([0,1])$ y un número $\lambda \in \mathbb{R}$, entonces la ecuación de resolvente es la siguiente
        \begin{equation*}            
            (T - \lambda Id)(x) = f
        \end{equation*}
        y buscamos encontrar si existe dicho $x$ que cumpla esta ecuación. Si existe, entonces $\lambda \in \rho(T)$, y si no existe, entonces $\lambda \in \sigma(T)$. Para resolver esta ecuación, basta con pasar todo a un solo lado y sacar factor común $x(t)$, quedando la ecuación de la siguiente forma
        \begin{equation*}
            (T - \lambda Id)(x)(t) = t \cdot x(t) - \lambda x(t) = (t - \lambda) x(t) = f(t) \iff x(t) = \frac{f(t)}{t - \lambda} \quad \forall t \in [0,1]
        \end{equation*}
        Ahora tenemos dos casos a estudiar
        \begin{itemize}
            \item Si $\lambda \notin [0,1]$, entonces $t - \lambda \neq 0$ para todo $t \in [0,1]$, por lo que $x(t) = \frac{f(t)}{t - \lambda}$ es una función continua (ya que $f$ es continua y el denominador no se anula), por lo que existe solución para cualquier $f$, y además es única, por lo que $\lambda \in \rho(T)$.
            \item Si $\lambda \in [0,1]$, entonces existirá un punto exacto en el dominio, al que llamamos $t_0 = \lambda$, donde el denominador se desvanece por completo: $(t_0 - \lambda) = 0$. Si intentamos evaluar la solución candidata en ese punto, tendríamos:$$x(\lambda) = \frac{f(\lambda)}{0}$$Si elegimos una función $f$ tal que $f(\lambda) \neq 0$, la división por cero genera una asíntota vertical (una discontinuidad que tiende a infinito). Es imposible encontrar una función continua $x(t)$ que satisfaga la igualdad en ese punto. Por lo tanto, no existe solución para esa función $f$, y por tanto $\lambda \in \sigma(T)$. Dado esto tenemos entonces que 
            \begin{equation*}
                [0,1]\subseteq \sigma(T)
            \end{equation*}
        \end{itemize}
        Juntando ambos cosas tenemos finalmente que 
        \begin{equation*}
            \sigma(T) = [0,1]
        \end{equation*}
        
        \item[$ii)$] Para determinar si $T$ es compacto o no, basta con estudiar si el espectro de $T$ sin el punto $0$ coincide con el conjunto de valores propios sin el punto $0$. En este caso vamos a demostrar que el conjunto de valores propios es vacío, es decir, $e(T) = \emptyset$, lo cual implica que $T$ no es compacto (de lo contrario se tendría $\sigma(T) \setminus \{0\} = e(T) \setminus \{0\}$). Para ello, basta con resolver la ecuación de valores propios 
        $$T(x) = \lambda x \iff T(x)(t) = \lambda x(t) \iff  t\cdot x(t) = \lambda x(t)\iff (t - \lambda) x(t) = 0 \quad \forall t\in [0,1]$$
        donde tenemos que $x$ debe ser una función no nula. Para que el producto de dos números reales sea cero, es obligatorio que el primero sea cero, el segundo sea cero, o ambos. Por lo tanto, en cada punto $t \in [0,1]$, la matemática nos obliga a elegir entre una de estas dos opciones:
        \begin{itemize}
            \item Opción A: el primer paréntesis se anula: $t = \lambda$.
            \item Opción B: la función se anula en ese instante: $x(t) = 0$.
        \end{itemize}
        Supongamos por un momento que sí existe un valor propio $\lambda$ con su correspondiente función propia $x(t) \neq 0$. Como hemos supuesto que $x \neq 0$, la definición nos garantiza que existe al menos un punto concreto, llamémoslo $t_0$, donde la función tiene una altura distinta de cero
        $$x(t_0) \neq 0$$
        Si en el punto $t_0$ la función se encuentra a una altura $x(t_0) \neq 0$, está obligada a mantenerse alejada del cero en todo un pequeño entorno (un intervalo abierto) alrededor de $t_0$, entonces
        $$x(t) \neq 0 \quad \forall t \in \text{Entorno}(t_0)$$
        lo cual es claramente un absurdo, pues $\lambda$ es una constante y no puede tener varios valores en un entorno de $t_0$. Por lo tanto, no existe ningún valor propio $\lambda$ y $e(T) = \emptyset$.
    \end{description}
    Antes de finalizar el ejercicio, calcularemos la norma del operador $T$.  \\ \\
    \noindent
    Sabemos que $\sigma(T) = [0,1] \subseteq [-\|T\|, \|T\|]$, lo cual implica que $\|T\| \geq 1$. Por otro lado, para cualquier función $x \in C^0([0,1])$, se cumple que
    \begin{equation*}
        T(x)(t) = t \cdot x(t) \implies \|T(x)\| = \sup_{t \in [0,1]} |t \cdot x(t)| \leq \sup_{t \in [0,1]} |x (t)| = \|x\|
    \end{equation*}
    lo cual implica que $\|T\| \leq 1$. Juntando ambas cosas, se concluye que $\|T\| = 1$.
\end{ejercicio}

\section{Operadores autoadjuntos}
\noindent
Sea $\mathcal{H}$ un espacio de Hilbert sobre $\mathbb{C}$, y $T \in BL(\mathcal{H})$, entonces dado $T$ autoadjunto, recordemos que tenemos que 
\begin{equation*}
    \|T\| = \sup_{\|u\|=1} |( T(u), u )| = \sup_{u \neq 0} \frac{|( T(u), u )|}{\|u\|^2}
\end{equation*}
veámos como podemos estimar el espectro.

\begin{teo}
    Sea $T \in BL(H)$, $H$ espacio de Hilbert  y $T$ autoadjunto, definimos
    \begin{equation*}
        m := \inf_{\substack{u \in H \\ \|u\|=1}} (T(u), u ) \quad \text{y} \quad M := \sup_{\substack{u \in H \\ \|u\|=1}} ( T(u), u )
    \end{equation*}
    tenemos entonces que 
    \begin{equation*}
        \sigma(T) \subseteq [m, M] \quad \text{y} \quad m, M \in \sigma(T)
    \end{equation*}
\begin{proof} % //TODO: leer demostración y ver si se copio bien
    % \underline{$1^{\text{er}}$ paso:} Si $\lambda < m \implies \lambda \in \rho(T)$. Dado $u \in H, \ u \neq 0$, se tiene $\langle T(u), u \rangle = \|u\|^2 \left\langle \frac{T(u)}{\|u\|}, \frac{u}{\|u\|} \right\rangle \geq m \|u\|^2$

    % \noindent $\left(\text{ya que } \left\| \frac{u}{\|u\|} \right\| = 1 \right)$.

    % \vspace{0.3cm}

    % \noindent $\langle T(u) - \lambda u, u \rangle = -\lambda \langle u, u \rangle + \langle T(u), u \rangle \geq \underbrace{(m - \lambda)}_{=: \alpha} \|u\|^2 = \alpha \|u\|^2, \quad \alpha > 0$


    % \noindent $a_\lambda(u, v) := \langle (T - \lambda I)(u), v \rangle$ es

    % \noindent una \textbf{forma bilineal coercitiva}: $a_\lambda(u, u) \geq \alpha \|u\|^2$

    % \vspace{0.4cm}

    % \noindent $\implies$ por el th. de Lax-Milgram: $\forall f \in H \ \exists! u \in H : \underbrace{a_\lambda(u, v)}_{\langle (T - \lambda I)(u), v \rangle} = \langle f, v \rangle \quad \forall v \in H$.

    % \vspace{0.4cm}

    % \noindent $\implies (T - \lambda I)(u) = f$ tiene una única soluz. $\forall f \in H \implies T - \lambda I$ es biyectiva. $\checkmark$ç

    % \noindent \underline{$2^{\text{o}}$ paso:} $\lambda = m \in \sigma(T)$. \hfill De hecho: dado $m$ fijado, consideremos la forma

    % \noindent bilineal $a_m(u, v) := \langle (T - mI)(u), v \rangle$ y mostremos que es

    % \noindent simétrica. $a_m(u, v) = \langle (T - mI)(u), v \rangle = \langle Tu, v \rangle - m \langle u, v \rangle = $
    % \[ = \langle u, T(v) \rangle - m \langle v, u \rangle = \langle T(v), u \rangle - m \langle v, u \rangle = \]
    % \[ = \langle T(v) - mv, u \rangle = a_m(v, u). \]

    % \vspace{0.3cm}

    % \noindent Además, $a_m(v, v) = \langle (T - mI)(v), v \rangle \geq 0$ \quad por def. de $m = \inf\dots$

    % \vspace{0.3cm}

    % \noindent $\leadsto a_m$ es casi un producto escalar $\leadsto$ llamémoslo $\langle \cdot, \cdot \rangle$.
    % \begin{center}
    %     $\downarrow$ \\
    %     Si $\langle u, u \rangle = 0$, no está dicho que $u = 0$.
    % \end{center}
    % \begin{center}
    %     $\downarrow$ \\
    %     \textit{pero esta propiedad no se necesita para demostrar} \\
    %     \textit{la desig. de Cauchy-Schwarz}
    % \end{center}
    % \noindent por lo tanto, para $\langle \cdot, \cdot \rangle$ vale Cauchy-Sch: $|\langle u, v \rangle| \leq \sqrt{\langle u, u \rangle} \sqrt{\langle v, v \rangle}$

    % \vspace{0.3cm}

    % \noindent $\implies |a_m(u, v)| = |\langle (T - mI)(u), v \rangle| \leq \sqrt{\langle (T - mI)(u), u \rangle} \sqrt{\langle (T - mI)(v), v \rangle}$

    % \noindent $\forall u, v \in H$.

    % \vspace{0.4cm}

    % \noindent Aplicamos la desigualdad con $v = (T - mI)(u) =: Au$
    % \begin{center}
    %     $\uparrow$ \\
    %     $A$
    % \end{center}

    % \noindent $\|A(u)\|^2 \leq \langle A(u), u \rangle^{\frac{1}{2}} \langle A(A(u)), A(u) \rangle^{\frac{1}{2}} \leq \langle A u, u \rangle^{\frac{1}{2}} \left( \|A\| \cdot \|A(u)\|^2 \right)^{\frac{1}{2}}$

    % \noindent $\implies \|A(u)\| \leq \langle A u, u \rangle^{\frac{1}{2}} \underbrace{\|A\|^{\frac{1}{2}}}_{c \text{ cost.}} = c \langle A(u), u \rangle^{\frac{1}{2}}$


    % \noindent $\implies \|(T - mI)(u)\| \leq c \langle (T - mI)(u), u \rangle^{\frac{1}{2}} \quad \forall u \in H$

    % \vspace{0.3cm}

    % \noindent Por def. de $m$, $\exists \{u_n\} : \|u_n\| = 1 \quad \text{y} \quad \langle T(u_n), u_n \rangle \longrightarrow m$.

    % \vspace{0.3cm}

    % \noindent $\|(T - mI)(u_n)\| \leq c \langle (T - mI)(u_n), u_n \rangle^{\frac{1}{2}} = c \left( \underbrace{\langle T(u_n), u_n \rangle}_{\downarrow} - m \underbrace{\langle u_n, u_n \rangle}_{\substack{\downarrow \\ 1}} \right)^{\frac{1}{2}} = $
    % \begin{center}
    %     $\begin{array}{ccc} \quad & m & \quad \end{array}$ \\
    %     $\underbrace{\hspace{5cm}}_{\downarrow}$ \\
    %     $0$
    % \end{center}

    % \noindent $\implies m \in \sigma(T)$, porque si $m \notin \sigma(T)$

    % \noindent es decir, $m \in \rho(T)$, entonces se tendría $T - mI$ invertible

    % \noindent con inversa continua $\implies (T - mI)^{-1}(T - mI)(u_n) \longrightarrow 0$

    % \noindent $\iff u_n \longrightarrow 0 \quad \text{absurdo} \quad (\|u_n\| = 1)$.

    % \vspace{0.4cm}

    % \noindent Análogamente se muestra que si $M < \lambda \implies \lambda \in \rho(T)$ y $M \in \sigma(T)$.
\end{proof}
\end{teo}

\observacion{
    Sabemos que $m\in \sigma(T)$, de hecho si $m$ fuera un valor propio (\tbf{no nulo}), entonces \underline{$m = \min$}, veámoslo. Si $m$ es un valor propio, entonces existe un vector propio $\bar{u}$ tal que $T(\bar{u}) = m \bar{u}$ y $\|\bar{u}\| = 1$. Entonces, se tiene que
    \begin{equation*}
        m = ( T(\bar{u}), \bar{u} ) = ( m \bar{u}, \bar{u} ) = m ( \bar{u}, \bar{u} ) = m \cdot 1 = m
    \end{equation*}
    y por lo tanto, $m$ es el mínimo de la función $( T(u), u )$ sobre la esfera unitaria $\{u \in H : \|u\| = 1\}$, es decir 
    $$m = \min_{\|u\|=1} \langle T(u), u \rangle$$
    Analógamente, si $M$ es un valor propio (\tbf{no nulo}), entonces \underline{$M = \max$}.
}

\begin{observacion}
    Es importante notar que los vectores propios no tienen porque ser unitarios, es decir, podemos obtener un vector propio $u\in \mathcal{H}$ con norma distinta de $1$, pero como el operador es lineal, entonces 
    $$T(u) = T\left(\|u\| \cdot \frac{u}{\|u\|}\right) = \|u\| \cdot T\left(\frac{u}{\|u\|}\right)$$
    por lo que $\frac{u}{\|u\|}$ también es un vector propio asociado al mismo valor propio. \\ \\
    \noindent
    Por lo tanto, sin pérdida de generalidad, podemos asumir que los vectores propios tienen norma $1$, ya que la clave radica en que las soluciones de la ecuación de autovalores no son vectores aislados, sino subespacios vectoriales completos (líneas, planos, etc.), por lo que no existe un vector propio único, sino que existe una infinidad de vectores propios asociados a un mismo valor propio, y entre ellos se encuentra uno con norma $1$.
\end{observacion}

\begin{coro}
    Sea $T$ autoadjunto, entonces tenemos que 
    \begin{equation*}
        \sigma(T)=\{0\}\implies T \equiv 0
    \end{equation*}
\begin{proof}
    Sea $m=M=0$, entonces $( T(u), u ) = 0 \quad \forall u \in H$. Por otro lado, se tiene que
    \begin{equation*}
        \|T\| = \sup_{\|u\|=1} |( T(u), u )| = 0
    \end{equation*}
    lo cual implica que $T \equiv 0$. \\
\end{proof}
\end{coro}

\noindent
Veámos que hemos obtenido, partiendo de $T\in BL(\mathcal{H})$ autoadjunto.
\begin{itemize}
    \item $\sigma(T)=\{0\} \iff T \equiv 0$.
    \item Sobre $\mathbb{C}$ los valores propios son todos reales ($\lambda \in \mathbb{R}$), al igual que para las matrices simétricas. De hecho, si $\lambda \neq 0$ es un valor propio con vector propio $u : \|u\| = 1 \implies T(u) = \lambda u$, entonces
    \begin{equation*}        
        \begin{cases}
            ( T(u), u ) = \lambda ( u, u ) = \lambda  \\
            ( T(u), u )=( u, T(u) ) = \bar{\lambda} ( u, u ) = \bar{\lambda}
        \end{cases}
        \Longrightarrow \lambda = \bar{\lambda} \iff \lambda \in \mathbb{R}
    \end{equation*}
    lo cual implica que $\lambda$ es real si $T$ es autoadjunto.
    \item Los espacios propios correspondientes a valores propios distintos son ortogonales entre sí, veámoslo. \\ 
    \noindent
    Sean $y_1$ tal que $T(y_1) = \lambda_1 y_1$, e $y_2$ tal que $T(y_2) = \lambda_2 y_2$ con $\lambda_1 \neq \lambda_2$. Supongamos $\lambda_1 \neq 0$, entonces
    \begin{equation*}
        y_1=\dfrac{T(y_1)}{\lambda_1} \implies ( y_1, y_2 ) =\left( \frac{T(y_1)}{\lambda_1}, y_2 \right) = \frac{1}{\lambda_1} \left( y_1, T(y_2) \right) = \frac{1}{\lambda_1} \left( y_1, \lambda_2 y_2 \right) = \frac{\lambda_2}{\lambda_1} \left( y_1, y_2 \right)
    \end{equation*}
    donde para que se de dicha igualdad debemos tener o $\lambda_1 = \lambda_2$ (lo cual es absurdo porque hemos supuesto que son distintos) o $( y_1, y_2 ) = 0$. Por lo tanto, los espacios propios correspondientes a valores propios distintos son ortogonales entre sí. \\
\end{itemize}

\noindent
Demostrado lo anterior, vamos a resumirlo en un resultado claro.

\begin{prop}
    Sea $T \in BL(\mathcal{H})$ autoadjunto, entonces:
    \begin{itemize}
        \item $\sigma(T) =\{0\}\iff T\equiv 0$.
        \item $\sigma(T) \subseteq \mathbb{R}$.
        \item Si $\lambda_1, \lambda_2 \in \sigma(T)$ con $\lambda_1 \neq \lambda_2$, entonces los espacios propios correspondientes a $\lambda_1$ y $\lambda_2$ son ortogonales entre sí. \\
    \end{itemize}
\end{prop}

\noindent
Ahora trabajaremos con operadores $T$ autoadjuntos y compactos no nulos ($T \neq 0$), veámos qué podemos decir sobre su espectro. Sabemos que \underline{existe al menos un valor propio $\lambda$ tal que $|\lambda|=\|T\|$}, veámoslo. \\ \\
\noindent
Sabemos que $m,M\in \sigma(T)$, y si son diferentes de cero, son valores propios, puesto que antes demostramos que 
$$\sigma(T)\setminus\{0\}=e(T)\setminus\{0\}$$
para dimensión finita de $\mathcal{H}$ y para dimensión infinita tenemos el apartado $b)$ del Teorema de la alternativa de Fredhold. Entonces, tenemos cuatro posibilidades:
\begin{itemize}
    \item $m = M = 0 \implies \sigma(T) = \{0\} \implies T \equiv 0$, no nos vale este caso.
    \item $m = 0, \ M > 0 \implies M$ es un valor propio y $M = \|T\|$ (por la caracterización de la norma de los op. autoadjuntos).
    \item $m < 0, \ M = 0 \implies m$ es un valor propio y $|m| = \|T\|$
    \item $m \neq 0, \ M \neq 0 \implies m$ y $M$ son valores propios y $\|T\| = \max\{|m|, |M|\}$.
\end{itemize}

\begin{teo}
    Sea $H$ espacio de Hilbert separable, $T \in BL(\mathcal{H})$ autoadjunto y compacto. Entonces existe una base hilbertiana, es decir, un sistema ortonormal completo formado por vectores propios de $T$.
\begin{proof} % //TODO: leer demostración y ver si se copio bien
    % Si $T \equiv 0$ entonces $\text{Ker}(T) = H$. \quad (obv)

    % \vspace{0.3cm}

    % \noindent Supongamos $T \neq 0 \implies \exists$ un valor propio no nulo. Sea $\{\lambda_n\}_{n \in \mathcal{J}}$ el conjunto de los valores propios distintos entre sí y diferentes de cero.

    % \vspace{0.3cm}

    % \noindent $\mathcal{J}$ o es finito ($\mathcal{J} = \{1, \dots, K\}$) o es numerable ($\mathcal{J} = \mathbb{N}^+$).
    % \begin{flushright}
    %     $_{\hookrightarrow \text{obs: si } \lambda_0 \text{ no es valor propio } E_0 = \{0\}.}$
    % \end{flushright}

    % \noindent Definimos $\lambda_0 := 0$, \quad $E_0 := \text{Ker}(T)$, \qquad $E_n := \text{Ker}(T - \lambda_n I)$

    % \vspace{0.3cm}

    % \noindent $0 \leq \dim E_0 \leq +\infty, \qquad 0 < \dim E_n < +\infty$

    % \vspace{0.3cm}

    % \noindent $E_n \perp E_m \qquad \forall n, m \geq 0$

    % \vspace{0.3cm}

    % \noindent Sea $F$ el espacio generado por los subespacios $E_n$ para $n \geq 0$.

    % \vspace{0.4cm}

    % \noindent $\bullet$ \textbf{Afirmo que \textcolor{red}{\underline{$F$ es denso en $H$}}.}

    % \vspace{0.4cm}

    % \noindent \underline{De hecho:} \quad (i) $T(F) \subseteq F \quad \leadsto$ por def. de $F$, si $x \in F$, $x = \sum_{\text{finita}} \alpha_i x_i$

    % \vspace{0.2cm}

    % \noindent $\implies T(x) = \sum_{\text{finita}} \alpha_i T(x_i) = \sum \alpha_i \lambda_i x_i \in F$


    % \noindent (ii) $T(F^\perp) \subseteq F^\perp$, porque si $u \in F^\perp$ y $v \in F$, entonces $\langle T(u), v \rangle = $
    % \[ = \langle u, T(v) \rangle = 0 \implies T(u) \in F^\perp. \]
    % \begin{center}
    %     $\begin{array}{cc} \uparrow & \uparrow \\ F^\perp & F \end{array}$
    % \end{center}

    % \noindent Sea $T_0 = T_{|_{F^\perp}} \quad \leadsto$ autoadjunto y compacto.

    % \vspace{0.3cm}

    % \noindent (iii) $\sigma(T_0) = \{0\} \quad \leadsto$ si hubiera un $\lambda \in \sigma(T_0) \setminus \{0\}$, entonces $\lambda$ sería un valor propio

    % \noindent $\implies \exists u \in F^\perp : u \neq 0$ y $T_0(u) = \lambda u$, es decir, $\lambda$ es un valor propio de $T$
    % \begin{center}
    %     $\begin{array}{c} \parallel \\ T(u) \end{array}$
    % \end{center}

    % \noindent es decir, un $\lambda_{\bar{n}}$ para un cierto $\bar{n} \implies u \in E_{\bar{n}} \subseteq F \implies u \in F \cap F^\perp \quad \text{absurdo}$

    % \noindent ($\implies u = 0$). Por lo tanto, $T_0 \equiv 0$ (porque su espectro es $\{0\}$), es decir,

    % \noindent $T \equiv 0$ sobre $F^\perp \implies F^\perp \subseteq \text{Ker}(T) = E_0 \subseteq F \implies F^\perp \subseteq F \implies F^\perp = \{0\}$.

    % \vspace{0.4cm}

    % \noindent Hemos por lo tanto mostrado que $F^\perp = \{0\}$, pero entonces $F$ es denso

    % \noindent en $H$, porque si $\overline{F} \neq H \implies F$ tendría un complemento ortogonal

    % \noindent (por el th. de descomposición ortogonal). $\checkmark$

    % \vspace{0.4cm}

    % \noindent Tomemos en cada $E_n$, para $n \geq 1$, una base ortonormal

    % \noindent (finita) y en $E_0$ una base ortonormal (a lo sumo numerable, porque

    % \noindent $H$ es separable). La unión de estas bases es un

    % \noindent \underline{sistema ortonormal numerable $\{e_k\}_{k \in \mathbb{N}}$} que es \underline{completo}.


    % \noindent porque el espacio generado por los $e_k$ es \underline{$F = \text{span}\{e_k\}$}, el cual es

    % \noindent \underline{denso en $H$}. De hecho, si no fuera completo $\exists y \in H$ t.q. $y \neq 0$ e

    % \noindent $y \perp e_k \quad \forall k \implies y \perp \text{span}\{e_k\} = F \implies y \in F^\perp \quad \text{absurdo}$.
\end{proof}
\end{teo}

\begin{observacion}
    El teorema anterior es un resultado muy importante, porque nos dice que en este caso, el operador $T$ se puede diagonalizar completamente, es decir, se puede representar como una matriz diagonal con los valores propios en la diagonal y los vectores propios como columnas. Esto es análogo a lo que ocurre con las matrices simétricas en espacios finitos, pero en este caso estamos trabajando con operadores en espacios de Hilbert de dimensión infinita.
\end{observacion}

\begin{teo}
    Sea $X$ un espacio de Banach sobre $\C$ y $T \in BL(X)$, entonces 
    $$\sigma(T) \neq \emptyset$$
\begin{proof}

    % \noindent \textbf{\textcolor{red}{\underline{Teorema:}}} $X$ espacio de Banach sobre $\mathbb{C}$ y $T \in BL(X)$. Entonces $\sigma(T) \neq \emptyset$.

    % \vspace{0.4cm}

    % \noindent \textbf{\textcolor{red}{\underline{dim:}}} Sean $\lambda, \mu \in \rho(T)$, con $\lambda \neq \mu$. Definimos $(T - \lambda I)^{-1} =: T(\lambda)$.

    % \vspace{0.3cm}

    % \noindent $T(\lambda) = T(\lambda)(T - \mu I)T(\mu) = T(\lambda) \big[ T - \lambda I + (\lambda - \mu)I \big] T(\mu) = $
    % \[ = \big[ I + (\lambda - \mu)T(\lambda) \big] T(\mu) \implies T(\lambda) - T(\mu) = (\lambda - \mu)T(\lambda)T(\mu) \qquad (*) \]

    % \vspace{0.4cm}

    % \noindent Sea $f \in [BL(X)]^*$ un funcional lineal $\quad \leadsto$ definimos $\quad \varphi(\lambda) := f(T(\lambda)) \quad \forall \lambda \in \rho(T)$.

    % \noindent $\varphi: \rho(T) \subseteq \mathbb{C} \longrightarrow \mathbb{C}$

    % \vspace{0.4cm}

    % \noindent Observamos que $|\varphi(\lambda)| = |f(T(\lambda))| \leq \|f\| \cdot \|T(\lambda)\|$.

    % \vspace{0.3cm}

    % \noindent $\frac{\varphi(\lambda) - \varphi(\mu)}{\lambda - \mu} = \frac{f(T(\lambda)) - f(T(\mu))}{\lambda - \mu} = \frac{f(T(\lambda) - T(\mu))}{\lambda - \mu} \overset{(*)}{=} $

    % \vspace{0.2cm}

    % \noindent $= \frac{(\lambda - \mu) f(T(\lambda) T(\mu))}{(\lambda - \mu)} = f(T(\lambda) T(\mu)) \xrightarrow{\lambda \to \mu} f(T(\mu)^2)$.

    % \vspace{0.4cm}

    % \noindent El límite existe, es decir, $\varphi$ es derivable en $\mathbb{C} \quad \forall \lambda \in \rho(T) \implies \varphi$ es holomorfa.


    % \noindent Si $|\lambda| > \|T\|$, entonces $\lambda \in \rho(T)$.

    % \vspace{0.3cm}

    % \noindent $T(\lambda) = (T - \lambda I)^{-1} = \left[ \lambda \left( \frac{T}{\lambda} - I \right) \right]^{-1} = \lambda^{-1} \left( \frac{T}{\lambda} - I \right)^{-1} = (-\lambda^{-1}) \left( I - \frac{T}{\lambda} \right)^{-1}$
    % \begin{flushright}
    %     $\parallel$ \\
    %     $\displaystyle\sum_{k=0}^{+\infty} \left( \frac{T}{\lambda} \right)^k$ \quad \textit{serie geom.}
    % \end{flushright}

    % \noindent $\left\| \frac{T}{\lambda} \right\| = \frac{\|T\|}{|\lambda|} < 1$

    % \vspace{0.2cm}

    % \noindent $\left\| \displaystyle\sum_{k=0}^{+\infty} \left( \frac{T}{\lambda} \right)^k \right\| \leq \displaystyle\sum_{k=0}^{+\infty} \left( \frac{\|T\|}{|\lambda|} \right)^k = \frac{1}{1 - \left\| \frac{T}{\lambda} \right\|}$

    % \vspace{0.4cm}

    % \noindent Por lo tanto, $\|T(\lambda)\| \leq \frac{1}{|\lambda|} \frac{1}{1 - \left\| \frac{T}{\lambda} \right\|} \implies \|T(\lambda)\| \longrightarrow 0 \quad \text{si } |\lambda| \to +\infty$

    % \vspace{0.4cm}

    % \noindent Por lo tanto, $\varphi(\lambda) = f(T(\lambda))$ es acotada, además de holomorfa, y $\varphi(\lambda) \longrightarrow 0$

    % \noindent para $|\lambda| \to +\infty \quad \forall f \in [BL(X)]^*$.

    % \vspace{0.4cm}

    % \noindent Si $\rho(T) = \mathbb{C}$, entonces $\varphi: \mathbb{C} \longrightarrow \mathbb{C}$ es holomorfa y acotada, y por lo tanto,

    % \noindent por Liouville, es constante. $\quad \begin{array}{c} \varphi(\lambda) \longrightarrow 0 \\ |\lambda| \to +\infty \end{array} \implies \varphi(\lambda) \equiv 0 \quad \forall f \in [BL(X)]^*$.

    % \vspace{0.2cm}

    % \begin{center}
    %     $\uparrow$ \\
    %     \textit{Hahn-Banach}
    % \end{center}
    % \noindent $\implies f(T(\lambda)) = 0 \quad \forall f \in [BL(X)]^* \implies T(\lambda) = 0 \quad \forall \lambda \in \mathbb{C} = \rho(T) \quad \text{absurdo}$

    % \vspace{0.2cm}

    % \noindent lo cual es absurdo porque $T(\lambda) = (T - \lambda I)^{-1}$.

    % \vspace{0.4cm}

    % \noindent Por lo tanto, $\rho(T) \neq \mathbb{C} \implies \sigma(T) \neq \emptyset$.
\end{proof}
\end{teo}

\observacion{   
    Sea $\mathcal{H}$ espacio de Hilbert de dimensión infinita, observamos que dependiendo de si $0$ es valor o no propio, tenemos una cantidad de valores propios u otra, véamoslo:
    \begin{itemize}
        \item Si $0$ no es un valor propio, entonces $e(T)$ es numerable.
        \item Si $0$ es un valor propio de multiplicidad finita, entonces $e(T)$ es numerable.
        \item Si $0$ es un valor propio de multiplicidad geométrica infinita, entonces $e(T)$ puede ser finito.
    \end{itemize}
}
\begin{prop}
    Todo operador autoadjunto y compacto $T$ se aproxima con operadores con imagen de dimensión finita.
\begin{proof}
    Primero sabemos que cualquier vector $u \in H$ se puede escribir como una serie de Fourier usando los vectores propios
    $$u = \sum_{n} ( u, e_n ) e_n$$
    donde $S=\{e_n\}$ es una base ortonormal, y si agrupamos todos los términos que corresponden al mismo valor propio $\lambda_n$, llamamos a ese bloque $u_n = ( u, e_n ) e_n$, vector propio. ya que como $e_n$ es un vector propio, este nuevo vector agrupado $u_n$ (vector propio por escalar $(u,e_n)$) sigue siendo un vector propio (porque un múltiplo de un vector propio es un vector propio). \\ \\
    \noindent
    Si aplicamos el operador $T$ a nuestro vector $u$, por linealidad el operador entra dentro del sumatorio y ataca a cada vector propio $u_n$
    $$T(u) = T\left(\sum_{n} u_n\right) = \sum_{n} T(u_n) = \sum_{n=1}^{+\infty} \lambda_n u_n$$
    donde guardamos el índice $n=0$ para el caso en que el cero sea un autovalor, $\lambda_0 = 0$, el cual desaparece de la suma porque multiplicar por cero anula el término. \\ \\
    \noindent
    Para demostrar que $T$ se puede aproximar por operadores de dimensión finita, haremos un ``corte'' (un truncamiento) en el término $k$ de la serie y define un nuevo operador llamado $T_k$ como
    $$T_k(u) := \sum_{n=1}^{k} \lambda_n u_n$$
    Ahora queremos demostrar que si hacemos el corte cada vez más lejos ($k \to \infty$), el operador truncado $T_k$ se pega tanto a $T$ que el error comete un límite hacia cero, veámoslo
    \begin{equation*}
        \|(T - T_k)(u)\| = \left\| \sum_{n=k+1}^{+\infty} \lambda_n u_n \right\| \leq \sup_{n \geq k+1} |\lambda_n| \cdot \left\| \sum_{n=k+1}^{+\infty} u_n \right\| \overset{(*)}{\longrightarrow} 0 \quad \text{si } k \to +\infty
    \end{equation*}
    donde en $(*)$ hemos usado que la serie original para el vector $u$ era perfectamente convergente en el espacio de Hilbert, esta parte no es más que el resto de una serie convergente. Por las propiedades fundamentales de las series, a medida que avanzas hacia el infinito, el resto de cualquier serie que converge tiene la obligación matemática de tender a cero.
\end{proof}
\end{prop}

\section{Formulación variacional}
\noindent
En esta sección trabajaremos con $\mathcal{H}$ un espacio de Hilbert separable, $T\in BL(\mathcal{H})$ un operador autoadjunto y compacto. \\ \\
\noindent
Al tener $T$ autoadjunto y compacto, sabemos que los valores propios son reales y por tanto tienen signo: o son mayores o iguales que cero ($\geq 0$), o son menores o iguales que cero ($\leq 0$). Gracias a que \underline{no se pueden acumular en ningún otro sitio que no sea el origen} (Teorema \ref{teo:fredholm}), la matemática nos da permiso para hacer algo imposible en otros conjuntos de números reales: podemos hacer una fila india ordenada con ellos. \\
\noindent
Vamos a hacer dos listas ordenadas:
\begin{itemize}
    \item La lista de los positivos ($\lambda_m^+$): Se ordenan de mayor a menor (decreciente) marchando hacia el cero: $\lambda_1^+ \geq \lambda_2^+ \geq \dots \geq 0$.
    \item La lista de los negativos ($\lambda_m^-$): Se ordenan de menor a mayor (creciente, es decir, de los más negativos a los que se acercan al cero): $\lambda_1^- \leq \lambda_2^- \leq \dots \leq 0$.
\end{itemize}
Un punto crítico a tener en cuenta es que \textit{repitiremos cada valor propio tantas veces como indica su multiplicidad geométrica}, es decir, si un valor propio $\lambda$ tiene multiplicidad geométrica $k$, entonces aparecerá $k$ veces en la lista, por ejemplo si el número $5$ tiene un espacio propio de dimensión 2 (hay dos vectores independientes que responden al 5), en nuestra lista lo escribiremos dos veces: $\lambda_1^+ = 5$ y $\lambda_2^+ = 5$. Esto es importante porque la multiplicidad geométrica nos indica cuántos vectores propios linealmente independientes hay asociados a ese valor propio, y al repetirlo en la lista, estamos reflejando esa información de manera explícita. \\ \\

\noindent
Tomando las sucesiones de vectores propios normalizados como sigue
\begin{equation*}
    \{u_n^+\} \text{ para los positivos}, \ \{u_n^-\} \text{ para los negativos}, \quad \|u_n^+\| = \|u_n^-\| = 1
\end{equation*}
podemos crear una base ortonormal en $\mathcal{H}$ como sigue
\begin{equation*}
    \{u_m^+, u_n^-\} \text{ es una base ortonormal en } \mathcal{H}
\end{equation*}
Finalmente definiremos los siguientes subespacios vectoriales
\begin{equation*}
    V_m^+ = \text{span}\{u_1^+, \dots, u_m^+\} \quad \text{y} \quad V_n^- = \text{span}\{u_1^-, \dots, u_n^-\}, \quad V_0 = \{0\}
\end{equation*}
    
\begin{teo}
    Dado $\mathcal{H}$ espacio de Hilbert separable, $T \in BL(\mathcal{H})$ autoadjunto y compacto, entonces 
    \begin{equation*}
        \lambda_n^+ = ( T(u_n^+), u_n^+ ) = \max_{\substack{v \neq 0 \\ v \perp V_{n-1}}} \frac{( T(v), v )}{\|v\|^2} = \min_{\substack{V \subseteq H \\ \dim V = n-1}} \left( \max_{\substack{v \neq 0 \\ v \perp V}} \frac{( T(v), v )}{\|v\|^2} \right) \quad \forall n\geq 1
    \end{equation*}
    En particular tenemos
    \begin{equation*}
        \lambda_1^+ = \max_{\substack{v \in H \\ v \neq 0}} \frac{( T(v), v )}{\|v\|^2}
    \end{equation*}
    Análogamente para los negatios tenemos
    \begin{equation*}
        \lambda_n^- = ( T(u_n^-), u_n^- ) = \min_{\substack{v \neq 0 \\ v \perp V_{n-1}}} \frac{( T(v), v )}{\|v\|^2} = \max_{\substack{V \subseteq H \\ \dim V = n-1}} \left( \min_{\substack{v \neq 0 \\ v \perp V}} \frac{( T(v), v )}{\|v\|^2} \right) \quad \forall n\geq 1
    \end{equation*}
    y en particular tenemos
    \begin{equation*}
        \lambda_1^- = \min_{\substack{v \in H \\ v \neq 0}} \frac{( T(v), v )}{\|v\|^2}
    \end{equation*}
\begin{proof}
    Demostraremos solo la primera parte. Sea $u_n^+$ el vector propio relativo a $\lambda_n^+$ con $\|u_n^+\| = 1$, entonces
    \begin{equation*}
        T(u_n^+) = \lambda_n^+ u_n^+ \implies ( T(u_n^+), u_n^+ ) = \lambda_n^+ ( u_n^+, u_n^+ ) \overset{(*)}{=} \lambda_n^+
    \end{equation*}
    donde en $(*)$ hemos usado que $\|u_n^+\|=(u_n^+,u_n^+) = 1$. 
\end{proof}
\end{teo}

\begin{ejercicio}
    Sean $X,Y$ espacios de Banach, $T$ un operador lineal tal que $T: X \longrightarrow Y$, y $\{f_i\}_{i \in \mathcal{I}} \subseteq Y^*$ una familia que ``separa los puntos de $Y$'', es decir
    \begin{equation*}
        \forall y_1 \neq y_2, \ \exists j \in \mathcal{I} \text{ tal que } f_j(y_1) \neq f_j(y_2)
    \end{equation*}
    Demostrar que si $f_i \circ T$ es continuo $\forall i \in \mathcal{I}$, entonces $T$ es continuo. \\ \\
    \noindent
    Vamos a usar el Teorema del Gráfico Cerrado, para ello entonces debemos demostrar que $T$ es cerrada, lo que equivale a demostrar que el gráfico $\Gamma_T$ es cerrado, es decir
    \begin{equation*}
        \forall (x_n) \subseteq X\ : \ x_n \longrightarrow x \text{ y } T(x_n) \longrightarrow y \implies y = T(x)
    \end{equation*}
    Supondremos las hipótesis, es decir, una sucesión cualquier $(x_n)$ que cumple
    \begin{equation*}
        x_n \longrightarrow x \text{ y } T(x_n) \longrightarrow y
    \end{equation*}
    entonces
    \begin{equation*}
        \lim_{n \to \infty} f_i(T(x_n)) \overset{(*)}{=} f_i\left(\lim_{n \to \infty} T(x_n)\right) = f_i(y)
    \end{equation*}
    donde en $(*)$ hemos usado que $f_i$ continua por pertenecer a $Y^*$. Mirando el Diagrama \ref{fig:diagrama-ejercicio1}, sabemos que $f_i \circ T$ es una aplicación de $X$ en $\mathbb{R}$ (y continua por hipótesis), y recordando que $x_n \longrightarrow x$, entonces
    $$\lim_{n \to \infty} (f_i \circ T)(x_n) = (f_i \circ T)(\lim_{n \to \infty} x_n) = (f_i\circ T)(x)=f_i(T(x))$$
    Uniendo ambos resultados tenemos que 
    \begin{equation*}
        f_i(T(x)) = f_i(y) \quad \forall i \in \mathcal{I}
    \end{equation*}
    Lo que queda ahora es pensar que aún siendo las imágenes iguales, ¿es posible que $y$ y $T(x)$ sean vectores distintos? Esto sabemos que no es posible pues en el caso de $y\neq T(x)$, al tener la propiedad de que la familia ``separa los puntos de $Y$'', entonces tendríamos 
    \begin{equation*}
        f_j(y) \neq f_j(T(x)) \quad \text{para algún } j \in \mathcal{I}
    \end{equation*}
    lo que es un absurdo con lo obtenido. Queda por tanto demostrado que
    \begin{equation*}
        y = T(x)
    \end{equation*}

\begin{figure}
    \centering
    % Definimos el diagrama con TikZ
    \begin{tikzpicture}[baseline=(X.base), every node/.style={on grid}, node distance=2cm]
        % Posicionamiento de los elementos en línea recta
        \node (X) at (0,0) {$X$};
        \node (Y) at (2,0) {$Y$};
        \node (R) at (4,0) {$\mathbb{R}$};
        
        % Flechas superiores rectas
        \draw[->] (X) -- node[above] {$T$} (Y);
        \draw[->] (Y) -- node[above] {$f_i$} (R);
        
        % Flecha inferior curva (semicírculo por abajo) desde X hasta R
        % "bend right=50" controla la curvatura hacia abajo
        \draw[->] (X) to[bend right=50] node[below=2pt] {$f_i \circ T$} (R);
    \end{tikzpicture}
    \caption{Diagrama de la aplicación $f_i \circ T$.}
    \label{fig:diagrama-ejercicio1}
\end{figure}

\end{ejercicio}


\begin{teo}[Teorema  Ascoli-Arzelà]
    Sea $K \subset \mathbb{R}^N$ un conjunto compacto, y sea $C(K)$ el espacio de las funciones continuas sobre $K$ equipado con la norma del supremo ($\|\cdot\|_\infty$). Un subconjunto de funciones $A \subset C(K)$ es relativamente compacto  si y solo si cumple simultáneamente dos condiciones:
    \begin{itemize}
        \item \tbf{Acotación uniforme:} El conjunto de funciones no se escapa hacia el infinito verticalmente, es decir
        \begin{equation*}
            \exists M > 0 \text{ tal que } \|f\|_\infty \leq M \quad \forall f \in A
        \end{equation*}
        \item \tbf{Equicontinuidad uniforme:} Todas las funciones del conjunto comparten el mismo ``ritmo'' de continuidad. Es decir, para cualquier $\varepsilon > 0$, existe un mismo $\delta > 0$ tal que
        \begin{equation*}
            \text{Si } |x - y| < \delta \implies |f(x) - f(y)| < \varepsilon \quad \forall x,y \in K, \ \forall f \in A
        \end{equation*}
    \end{itemize}
\end{teo}

\begin{ejercicio}
    Sea $X=L^2((0,1))$ el espacio de Hilbert de las funciones de cuadro integrable, se define el \tbf{operador de Volterra} $T$ como sigue
    \Func{T}{X}{X}{u}{\int_{0}^{x} u(t) \, dt}
    Demuestra los siguientes puntos:
    \begin{itemize}
        \item[i)] $T$ es compacto.
        \item[ii)] $0 \in \sigma(T)$.
        \item[iii)] ¿$0 \in e(T)$?
    \end{itemize}
    Veámos la demostración de los puntos por separado.
    \begin{description}
        \item[$i)$] Para demostrar que $T$ es compacto, primero observamos que $T(u)$ es una función continua en el intervalo $[0,1]$, ya que la integral de una función de $L^2$ es continua. \\ \\
        \noindent
        Luego, si tomamos un conjunto acotado $B$ en $L^2((0,1))$, la imagen de este conjunto bajo $T$, es decir, $T(B)$, también será acotada en el espacio de funciones continuas $C([0,1])$. Recordemos que esto significa lo siguiente
        \begin{equation} \label{eq:acotacion-uniforme}
            \exists M > 0 \text{ tal que } \|u\|_{L^2} \leq M \quad \forall u \in B
        \end{equation}
        Vamos a hacer ahora dos demostraciones distintas
        \begin{itemize}
            \item \tbf{Acotación uniforme:} recordemos primero como funciona la norma infinito para $u\in B$ arbitrario
            \begin{equation*}
                \|T(u)\|_ \infty = \max_{x \in [0,1]} |T(u(x))|
            \end{equation*}
            entonces veamos como estimar el valor absoluto de la función para un $x\in [0,1]$ arbitrario
            \begin{align*}
                |T(u(x))| &\leq \int_{0}^{x} |u(t)| \cdot 1\, dt \overset{(*)}{\leq}  \left( \int_{0}^{x} |u(t)|^2 \, dt \right)^{1/2} \cdot \left( \int_{0}^{x} 1^2 \, dt \right)^{1/2}= \\
                &= \left( \int_{0}^{x} |u(t)|^2 \, dt \right)^{1/2} \cdot |x|^{1/2} \overset{(**)}{\leq} \|u\|_{L^2} \cdot |x|^{1/2} \leq \|u\|_{L^2}
            \end{align*}
            donde en $(*)$ hemos usado la desigualdad de Hölder, y en $(**)$ es, por definición, menor o igual que la norma completa de la función en todo el intervalo. Tenemos aplicado esto a todo $x\in [0,1]$ lo siguiente
            \begin{equation*}
                \|T(u)\|_\infty = \max_{x \in [0,1]} |T(u(x))| \leq \max_{x \in [0,1]} \|u\|_{L^2} = \|u\|_{L^2} \quad \forall u \in B
            \end{equation*}
            y recordando la ecuacion \ref{eq:acotacion-uniforme}, entonces tenemos 
            \begin{equation*}
                \|T(u)\|_\infty \leq \|u\|_{L^2}\leq M \quad \forall u \in B
            \end{equation*}
            y se tiene la acotación uniforme, es decir, $T(B)$ está acotado en $C([0,1])$.
            \item \tbf{Equicontinuidad uniforme:} para demostrar la equicontinuidad uniforme, debemos mostrar primero lo siguiente
            $$|T(u)(x) - T(u)(y)| \leq |x - y|^{1/2} \|u\|_{L^2} \leq M|x - y|^{1/2} \quad \forall u \in B, \forall x,y \in [0,1]$$
            para ellos usaremos lo mismo que en el caso anterior con la desiguald de Hölder, veámoslo
            \begin{equation*}
                |T(u)(x) - T(u)(y)| = \left| \int_{0}^{x} u(t) \, dt - \int_{0}^{y} u(t) \, dt \right| = \left| \int_{y}^{x} u(t) \, dt \right| \leq |x - y|^{1/2} \|u\|_{L^2}
            \end{equation*}
            donde tenemos por tanto demostrado la equicontinuidad uniforme.
        \end{itemize}
        Habiendo demostrado estos dos puntos, por el Teorema Ascoli-Arzelà, tenemos $\overline{T(B)}$ compacto en $C([0,1])$, y por tanto $T(B)$ es compacto en $L^2((0,1))$, ya que la norma de $L^2$ es más débil que la norma del supremo, es decir, si algo se pega mucho en la norma del supremo, también se va a pegar obligatoriamente en la norma de la integral al cuadrado. Por lo tanto, $T$ es compacto.

        \item[$ii)$] Para demostrar que $0 \in \sigma(T)$, basta con recordar que $T$ es compacto, y commo $X$ es de dimensión infinita, tenemos que 
        $$0 \in \sigma(T)$$
        Esto que hemos aplicado se demostró anteriormente en un ejericicio.
        \item[$iii)$] Para que $\lambda = 0$ sea un valor propio, tendríamos que ser capaces de encontrar una función propia $u \neq 0$ que al meterla en el operador diera cero
        $$T(u) = 0 \cdot u = 0$$
        En otras palabras, queremos ver si el núcleo del operador ($\text{Ker}(T)$) esconde funciones distintas de cero, para ello viendo la definición de dicho núcleo
        \begin{equation*}
            Ker(T)=\{u \in L^2((0,1)) : T(u) = 0\} = \{u \in L^2((0,1)) : T(u)(x) = \int_{0}^{x} u(t) \, dt = 0 \quad \forall x \in [0,1]\}
        \end{equation*}
        Para poder trabajar en el intervalo $[0,1]$, vamos a usar la función característica de un intervalo $[a,b] \subset [0,1]$, es decir, la función $\chi_{[a,b]}$ que vale $1$ en el intervalo $[a,b]$ y vale $0$ fuera de él. Entonces, si $u \in \text{Ker}(T)$, entonces
        \begin{equation*}
            \int_{0}^{1} u(t) \cdot \chi_{[0,x]}(t) \, dt = \int_{0}^{x} u(t) \, dt = 0 \quad \forall x \in [0,1]
        \end{equation*}
        Si la integral da cero para cualquier caja que empiece en $0$ ($\chi_{[0,x]}$), por resta de intervalos también tiene que dar cero para cualquier caja genérica que empiece en $a$ y termine en $b$ ($\chi_{[a,b]}$). Por ejemplo, si integras entre $0$ y $b$ y te da cero, y entre $0$ y $a$ y te da cero, la integral en el trozo intermedio $[a,b]$ es obligatoriamente cero, entonces tenemos que
        $$ \int_{0}^{1} u \cdot f \, dt = 0 \quad \forall f \in L^2((0,1))$$
        porque las funciones características son densas en $L^2$. Aquí está el corazón de la demostración. En la teoría de la medida de $L^2$, sabemos que cualquier función salvaje o continua $f \in L^2$ se puede aproximar tanto como queramos mediante combinaciones lineales de funciones de cajas (llamadas funciones escalonadas o simples). \\ \\
        \noindent
        Tenemos por tanto
        $$( u, f )_{L^2} = 0 \quad \forall f \in L^2 \implies u = 0$$
        es decir, el único vector propio asociado al valor propio $0$ es el vector nulo, por lo tanto $0 \notin e(T)$, ya que supusimos que podía ser no nulo y llegamos a la conclusión de que tenía que ser nulo, lo que es un absurdo. Por lo tanto, $0$ no es un valor propio de $T$.
    \end{description}
\end{ejercicio}

\subsection{Análisis espectral de ecuaciones en derivadas parciales (EDPs)}

\noindent
Cuando me pregunto si un operador tiene valores propios, quiero resolver la ecuación $-\Delta w = \lambda w$ en un dominio $\Omega$ con $w=0$ en el borde $\partial\Omega$, lo que viene a ser
\begin{equation*}
    \begin{cases} -\Delta w = \lambda w & \text{en } \Omega \\ w = 0 & \text{en } \partial\Omega \end{cases} \qquad \Omega \text{ abierto acotado, regular en } \mathbb{R}^N
\end{equation*}
donde si lo queremos resolver de la \underline{manera débil} sobre $H_0^1(\Omega)$, tendríamos que
\begin{equation*}
    \int_{\Omega} \nabla w \cdot \nabla \varphi \, dx = \lambda \int_{\Omega} w \cdot \varphi \, dx \qquad \forall \varphi \in H_0^1(\Omega)
\end{equation*}

\begin{teo}
    Existe una sucesión $\{w_m\} \subseteq H_0^1(\Omega)$ de funciones ortonormales completa en $L^2(\Omega)$, y existe una sucesión $\{\lambda_m\}$ de números reales positivos tal que $\lambda_m \longrightarrow +\infty$ y
    \begin{equation*}
        0 < \lambda_1 < \lambda_2 \leq \lambda_3 \leq \dots \quad \text{ tal que } \begin{cases} -\Delta w_m = \lambda_m w_m \\ \\ w_m \in H_0^1(\Omega) \cap C^\infty(\Omega) \end{cases}
    \end{equation*}
\begin{proof} % //TODO: leer demostración y ver que se copio bien

    % \noindent \textbf{\textcolor{red}{\underline{Dim:}}} $\forall f \in L^2(\Omega) \ \exists! \ u \in H_0^1(\Omega)$ t.q. $\int_{\Omega} \nabla u \cdot \nabla \varphi \, dx = \int_{\Omega} f \varphi \, dx \quad \forall \varphi \in H_0^1(\Omega)$

    % \vspace{0.3cm}

    % \noindent $\iff$ resuelve en forma débil $\begin{cases} -\Delta u = f & \text{en } \Omega \\ u = 0 & \text{en } \partial\Omega \end{cases}$

    % \vspace{0.4cm}

    % \noindent Por el th. de Lax-Milgram en $H_0^1(\Omega)$ aplicado a la forma bilineal

    % \noindent $a(u, v) = \int_{\Omega} \nabla u \cdot \nabla v \, dx, \quad u, v \in H_0^1(\Omega)$ \quad continua y coerciva

    % \vspace{0.2cm}

    % \noindent (por desig. de Poincaré: \quad $a(v, v) = \int_{\Omega} |\nabla v|^2 \, dx \geq k \|v\|_{H_0^1(\Omega)}^2$):

    % \vspace{0.3cm}

    % \noindent $\forall f \in L^2 \subseteq \big(H_0^1(\Omega)\big)^* = H^{-1}(\Omega) \quad \exists! \ u$ t.q. $a(u, v) = $

    % \noindent $= \int_{\Omega} f v \, dx \quad \forall v \in H_0^1(\Omega) \qquad \left( \text{es decir, soluz. en forma débil} \right)$

    % \vspace{0.5cm}

    % \noindent $G(f) \in H_0^1(\Omega)$ \quad única solución

    % \vspace{0.2cm}

    % \noindent $\int_{\Omega} \nabla (G(f)) \cdot \nabla v \, dx = \int_{\Omega} f v \, dx \quad \forall v \in H_0^1(\Omega) \qquad L^2(\Omega) \xrightarrow{\ G \ } H_0^1(\Omega)$

    % \vspace{0.5cm}

    % \noindent $\Omega$ acotado $\implies$ vale la inmersión de Sobolev: la inclusión $H_0^1(\Omega) \xhookrightarrow{\ i \ } L^2(\Omega)$

    % \noindent es compacta.

    % \begin{center}
    %     \begin{tikzpicture}[baseline=(H.base), every node/.style={on grid}, node distance=2.5cm]
    %         % Nodos del diagrama
    %         \node (L1) at (0,0) {$L^2(\Omega)$};
    %         \node (H) at (3,0) {$H_0^1(\Omega)$};
    %         \node (L2) at (6,0) {$L^2(\Omega)$};
            
    %         % Flechas superiores
    %         \draw[->] (L1) -- node[above] {$G$} (H);
    %         \draw[->] (H) -- node[above] {$i$} node[below=2pt] {\small cpt} (L2);
            
    %         % Flecha inferior curva (semicírculo por abajo)
    %         \draw[->] (L1) to[bend right=45] node[below=4pt] {$\tilde{G} := i \circ G \quad \text{compacta} \quad (G \text{ continua, } i \text{ cpt})$} (L2);
    %     \end{tikzpicture}
    % \end{center}

    % \vspace{0.4cm}

    % \noindent \underline{De hecho, $G$ es continua:} \quad $\|G(f)\|_{H_0^1(\Omega)}^2 = \int_{\Omega} |\nabla G(f)|^2 \, dx = \int_{\Omega} f \cdot u \, dx \leq$
    % \begin{flushright}
    %     $_{u = G(f)}$
    % \end{flushright}

    % \noindent $\leq \|f\|_{L^2} \|u\|_{L^2} \leq k \|f\|_{L^2} \|u\|_{H_0^1(\Omega)} \implies G \text{ acotado} \implies \text{continuo}$.

    % \vspace{0.6cm}

    % \noindent \underline{Mostremos que $\tilde{G}$ es autoadjunto}; es decir, mostremos que:
    % \[ \int_{\Omega} \tilde{G}(f) \cdot h \, dx = \int_{\Omega} f \cdot \tilde{G}(h) \, dx \qquad \forall f, h \in L^2(\Omega) \]

    % \noindent \left( \text{Escribo la eq. tanto para $f$ como para $h$:} \quad \int_{\Omega} \nabla (\tilde{G}(f)) \cdot \nabla (\tilde{G}(h)) \, dx = \int_{\Omega} f \cdot \tilde{G}(h) \, dx \right.
    % \begin{center}
    %     $\begin{array}{cc} \quad \ \ \uparrow & \qquad \qquad \qquad \ \ \uparrow \\ \text{\small $f$ soluz.} & \qquad \qquad \text{\small $h$ funz. test} \end{array}$
    % \end{center}

    % \noindent \left. \text{y} \quad \int_{\Omega} \nabla (\tilde{G}(h)) \cdot \nabla (\tilde{G}(f)) \, dx = \int_{\Omega} h \cdot \tilde{G}(f) \, dx \right)
    % \begin{center}
    %     $\begin{array}{cc} \quad \ \ \uparrow & \qquad \qquad \qquad \ \ \uparrow \\ \text{\small soluz.} & \qquad \qquad \text{\small test} \end{array}$
    % \end{center}

    % \vspace{0.3cm}

    % \noindent $\implies \tilde{G}$ es autoadjunto y compacto.

    % \vspace{0.3cm}

    % \noindent $\tilde{G}$ es def. positivo. \qquad $\int_{\Omega} \tilde{G}(f) \cdot f \, dx = \int_{\Omega} \nabla \tilde{G}(f) \cdot \nabla \tilde{G}(f) \, dx = \|\nabla \tilde{G}(f)\|_{H_0^1}^2 > 0 \quad \forall f \neq 0$.

    % \vspace{0.5cm}

    % \noindent Por resultados de teoría espectral de op. autoadjuntos y cpt, $\exists$ una base

    % \noindent ortonormal $\{w_m\} \subset L^2(\Omega)$ formada por funciones propias de $\tilde{G}$

    % \noindent correspondientes a $\{\mu_m\} : \mu_m \longrightarrow 0$ \quad (obs: $0$ no es valor propio

    % \noindent porque $G$ es inyectivo).

    % \vspace{0.4cm}

    % \noindent $\tilde{G}(w_m) = \mu_m w_m \iff w_m = \frac{1}{\mu_m} \tilde{G}(w_m), \qquad w_m \in H_0^1(\Omega)$

    % \vspace{0.3cm}

    % \noindent $\lambda_m = \frac{1}{\mu_m} \implies w_m = G(\lambda_m w_m) \implies$ por def. de $G$, $w_m$ es

    % \noindent la única soluz. de $\quad \int_{\Omega} \nabla w_m \cdot \nabla v \, dx = \int_{\Omega} \lambda_m w_m v \, dx \qquad \forall v \in H_0^1(\Omega)$.


\end{proof}
\end{teo}

\observacion{
    Sea $\{w_m\}$ una base ortonormal en $L^2$, pero nosotros queríamos una base ortonormal en $H_0^1(\Omega)$, entonces
    
    \begin{equation*}
         \varphi_m = \frac{1}{\sqrt{\lambda_m}} w_m \quad \text{ es ortonormal en } H_0^1(\Omega)
    \end{equation*}
}

\observacion{
    El valor propio $\lambda_1$ es simple (\underline{propiedad importante}), porque se puede caracterizar como:

    $$\lambda_1 = \min \frac{\int_{\Omega} |\nabla v|^2 \, dx}{\int_{\Omega} |v|^2 \, dx}$$
}
